
\section{Examples of brick polytopes}
\label{sec:interleaved-examples}

In this section, we give several examples of the brick polytopes realizing the corresponding subword complex. We use the normalization $\wt B(I^+)=0$ for the lifted vectors in the space  $M_Q$. We first give a root independent example $\mathcal P_1$ of type $\widetilde A_2$. Then we give a double root free but not root independent example $\mathcal P_2$ of type $A_3$. We have $\mathcal P_2\simeq \mathcal P_1\times [0,1]$ combinatorially, although the Coxeter groups and the words are totally different. Finally, we give a double root free but not root independent example of type $\widetilde A_2$.


\subsection{Example 1}
\label{ex:int-affine}

We take $W=\widetilde A_2$, $Q=s_0s_1s_0s_1s_2s_0s_2s_0 \text{ and }
\pi=\Dem(Q)=s_0s_1s_0s_2s_0.$
The length of $\pi$ is $5$ and $d=8-5=3$. Take $I^+=\{1,2,5\}$. We have
\[
R(I^+) =(\alpha_0,\alpha_1,2\alpha_0+\alpha_1+\alpha_2)
       =(\alpha_0,\alpha_1,\alpha_0+\delta),
\text{ where } \delta=\alpha_0+\alpha_1+\alpha_2.
\]
This example is root independent. In particular, it is double root
free, and $\Lambda_Q=\Z\alpha_0\oplus\Z\alpha_1\oplus\Z\alpha_2$.

Set
\[
M_Q=\Z U\oplus\Z V\oplus\Z W \text{ and } (U,V,W)=(e_1^{I^+},e_2^{I^+},e_5^{I^+}).
\]
The projection from $M_Q$ to $\Lambda_Q$ is
\begin{equation}\label{eq:int-affine-p}
p(u,v,w)=-u\alpha_0-v\alpha_1-w(2\alpha_0+\alpha_1+\alpha_2).
\end{equation}
It is an isomorphism of lattices in this example.

All facets, coordinates of lifted brick vectors after the normalization and lifted bases are listed as follows:

\begin{table}[htbp]
\centering
\caption{$\SC_{\widetilde A_2}(s_0s_1s_0s_1s_2s_0s_2s_0,s_0s_1s_0s_2s_0)$}
\label{tab:int-base}
\medskip
\begin{tabular}{@{}ccc@{}}
\toprule
$K$ & $z_K=(u,v,w)=\wt B(K)$ & $(e_i^K)_{i\in K}$\\
\midrule
$\{1,2,5\}$ & $(0,0,0)$ & $(U,V,W)$\\
$\{1,2,8\}$ & $(0,0,2)$ & $(U,V,-W)$\\
$\{1,5,6\}$ & $(0,2,0)$ & $(U,V+W,-V)$\\
$\{1,6,7\}$ & $(0,3,1)$ & $(U,W,-V-W)$\\
$\{1,7,8\}$ & $(0,3,2)$ & $(U,-V,-W)$\\
$\{2,3,5\}$ & $(1,0,0)$ & $(U+V,-U,W)$\\
$\{2,3,8\}$ & $(1,0,2)$ & $(U+V,-U,-W)$\\
$\{3,4,5\}$ & $(2,1,0)$ & $(V,-U-V,W)$\\
$\{3,4,8\}$ & $(2,1,2)$ & $(V,-U-V,-W)$\\
$\{4,5,6\}$ & $(2,2,0)$ & $(-U,V+W,-V)$\\
$\{4,6,7\}$ & $(2,3,1)$ & $(-U,W,-V-W)$\\
$\{4,7,8\}$ & $(2,3,2)$ & $(-U,-V,-W)$\\
\bottomrule
\end{tabular}
\end{table}


For the facet $I^+$, every $w(I^+,k)$ equals $\omega_{q_k}$ except $w(I^+,6)=\omega_0-\alpha_0 \text{ and } w(I^+,8)=\omega_0-\alpha_0-\alpha_1$. The brick vector $B(I^+)$ is $4\omega_0+2\omega_1+2\omega_2-2\alpha_0-\alpha_1$.
All brick vectors are given explicitly by
\begin{equation}\label{eq:int-affine-B}
B(K)=B(I^+)-u\alpha_0-v\alpha_1-w(2\alpha_0+\alpha_1+\alpha_2),
\text{ where } (u,v,w)=z_K.
\end{equation}
For example, $B(\{1,2,8\})=B(I^+)-2(2\alpha_0+\alpha_1+\alpha_2)$.

The convex hull of the lifted vectors is
\begin{equation}\label{eq:int-R}
\mathcal P_1=
\{(u,v,w)\mid 0\leq u\leq2,\ 0\leq v\leq3,\ 0\leq w\leq2,\
u-v\leq1,\ v-w\leq2\}.
\end{equation}


\begin{figure}[htbp]
\centering

\begin{tikzpicture}[
    x={(2.2cm,0cm)},
    y={(0.95cm,0.55cm)},
    z={(0cm,1.8cm)},
    line join=round,
    line cap=round,
    facetfill/.style={fill=blue!6,draw=none},
    visedge/.style={draw=black,thick},
    hidedge/.style={draw=black!65,thick,dashed},
    vtx/.style={circle,fill=black,inner sep=1.2pt},
    vlab/.style={font=\scriptsize\bfseries,text=blue!65!black,fill=white,inner sep=1.1pt},
    flab/.style={font=\scriptsize\bfseries,text=red!70!black,fill=white,inner sep=1.5pt,rounded corners=1pt}
]

%------------------------------------------------
% Vertices of P_1 = R
%------------------------------------------------
\coordinate (v125) at (0,0,0);
\coordinate (v128) at (0,0,2);
\coordinate (v156) at (0,2,0);
\coordinate (v167) at (0,3,1);
\coordinate (v178) at (0,3,2);

\coordinate (v235) at (1,0,0);
\coordinate (v238) at (1,0,2);

\coordinate (v345) at (2,1,0);
\coordinate (v348) at (2,1,2);

\coordinate (v456) at (2,2,0);
\coordinate (v467) at (2,3,1);
\coordinate (v478) at (2,3,2);

%------------------------------------------------
% Filled faces (back to front)
%------------------------------------------------
\fill[facetfill] (v345)--(v348)--(v478)--(v467)--(v456)--cycle; % r=4
\fill[facetfill] (v235)--(v238)--(v348)--(v345)--cycle;         % r=3
\fill[facetfill] (v128)--(v178)--(v478)--(v348)--(v238)--cycle; % r=8
\fill[facetfill] (v167)--(v178)--(v478)--(v467)--cycle;         % r=7
\fill[facetfill] (v156)--(v167)--(v467)--(v456)--cycle;         % r=6
\fill[facetfill] (v125)--(v128)--(v238)--(v235)--cycle;         % r=2
\fill[facetfill] (v125)--(v235)--(v345)--(v456)--(v156)--cycle; % r=5
\fill[facetfill] (v125)--(v128)--(v178)--(v167)--(v156)--cycle; % r=1

%------------------------------------------------
% Hidden edges (chosen viewing convention)
%------------------------------------------------
\draw[hidedge] (v128)--(v178);
\draw[hidedge] (v178)--(v167);
\draw[hidedge] (v167)--(v156);
\draw[hidedge] (v167)--(v467);

%------------------------------------------------
% Visible edges
%------------------------------------------------
\draw[visedge] (v125)--(v128);
\draw[visedge] (v128)--(v238);
\draw[visedge] (v238)--(v348)--(v478)--(v467)--(v456)--(v345)--(v235)--(v125);
\draw[visedge] (v125)--(v156);
\draw[visedge] (v156)--(v456);
\draw[visedge] (v235)--(v238);
\draw[visedge] (v235)--(v345);
\draw[visedge] (v345)--(v348);
\draw[visedge] (v456)--(v467);
\draw[visedge] (v467)--(v478);
\draw[visedge] (v178)--(v478);

%------------------------------------------------
% Vertices
%------------------------------------------------
\foreach \P in {125,128,156,167,178,235,238,345,348,456,467,478}
  \node[vtx] at (v\P) {};

%------------------------------------------------
% Vertex labels
%------------------------------------------------
\node[vlab,below left=2pt]  at (v125) {$125$};
\node[vlab,left=3pt]        at (v128) {$128$};
\node[vlab,left=3pt]        at (v156) {$156$};

% main correction here:
\node[vlab,right=8pt]        at ($(v167)+(-0.03,0.00,0.00)$) {$167$};
\node[vlab,above=7pt]       at ($(v238)+(0.00,0.00,0.02)$) {$238$};

\node[vlab,above left=2pt]  at (v178) {$178$};
\node[vlab,below=3pt]       at (v235) {$235$};
\node[vlab,below right=2pt] at (v345) {$345$};
\node[vlab,above right=2pt] at ($(v348)+(0.10,0.02,0.02)$) {$348$};
\node[vlab,below right=2pt] at (v456) {$456$};
\node[vlab,right=4pt]       at ($(v467)+(0.06,0.00,0.00)$) {$467$};
\node[vlab,above right=2pt] at (v478) {$478$};

%------------------------------------------------
% Face labels
%------------------------------------------------
\node[flab] at (-0.28,1.45,1.02) {$1$};
\node[flab] at (0.43,-0.16,1.00) {$2$};
\node[flab] at (1.52,0.35,1.08) {$3$};
\node[flab] at (2.22,1.95,1.02) {$4$};
\node[flab] at (1.06,1.00,-0.08) {$5$};
\node[flab] at (1.13,2.52,0.32) {$6$};
\node[flab] at (1.26,2.95,1.32) {$7$};
\node[flab] at (1.08,1.58,2.12) {$8$};

\end{tikzpicture}

\caption{The polytope $\mathcal P_1$. 
A vertex label $ijk$ stands for the facet $\{i,j,k\}$, and a face label corresponds to a position index in the subword complex.}
\label{fig:P1-polytope}

\end{figure}


\subsection{Example 2}
\label{ex:int-small}

Now we take $W=A_3$, $Q=s_1s_2s_3s_2s_1s_3s_2s_1s_2s_3,\text{ and }
\pi=w_0=s_2s_3s_2s_1s_2s_3$.
We have $\ell(\pi)=6$, $d=4$, and $I^+=\{1,2,3,5\}$.
The root configuration $R(I^+)$ is
\[
(\alpha_1,\alpha_2,\alpha_3,\alpha_1+\alpha_2).
\]
We have
\[
M_Q=\Z A\oplus\Z U\oplus\Z V\oplus\Z W,\qquad
(A,U,V,W)=(e_1^{I^+},e_2^{I^+},e_3^{I^+},e_5^{I^+}).
\]
For coordinates $(a,u,v,w)$, the projection and its kernel are
\begin{equation}\label{eq:int-small-kernel}
\begin{aligned}
p(a,u,v,w)&=-(a+w)\alpha_1-(u+w)\alpha_2-v\alpha_3,\\
\ker p&=\Z(A+U-W).
\end{aligned}
\end{equation}



For each of the twelve index sets $K$ in Table~\ref{tab:int-base}, we use
a relabeling $\sigma:i\to i+1, i\in\{1,2,3,6,7,8\}$, $\sigma(5)=5$, $\sigma(4)=6$ and set
\[
I_{0,K}=\{1\}\cup\sigma(K),\qquad
I_{1,K}=\{10\}\cup\sigma(K).
\]
These are precisely all the facets of this subword complex. If $z_K=(u,v,w)$, their lifted
brick vectors are
\begin{equation}\label{eq:int-small-vertices}
\wt B(I_{0,K})=(0,u,v,w),\qquad
\wt B(I_{1,K})=(5+u-w,u,v,w).
\end{equation}
Thus we have all lifted brick vectors and the polytope $\mathcal P_2$ is given by
\begin{equation}\label{eq:int-small-polytope}
\mathcal P_2=\{(a,u,v,w)\mid (u,v,w)\in\mathcal P_1,\quad
0\leq a\leq5+u-w\}.
\end{equation}

\begin{figure}[htbp]
\centering
\begin{tikzpicture}[
  x={(1.55cm,0cm)},
  y={(-0.40cm,0.30cm)},
  z={(0cm,0.85cm)},
  line cap=round,
  line join=round,
  lowedge/.style={draw=green!60!black,thick},
  topedge/.style={draw=blue!65!black,thick},
  fiber/.style={draw=orange!75!black,thick},
  fiberthin/.style={draw=orange!65!black,thin},
  important/.style={draw=orange!90!black,very thick},
  vlabel/.style={
    font=\scriptsize\bfseries,
    text=green!60!black,
    fill=white,
    inner sep=1pt
  },
  flabel/.style={
    font=\scriptsize\bfseries,
    text=red!70!black,
    fill=white,
    inner sep=1.2pt,
    rounded corners=1pt
  }
]

%================================================

%================================================
\foreach \id/\u/\v/\w in {125/0/0/0,128/0/0/2,156/0/2/0,167/0/3/1,178/0/3/2,235/1/0/0,238/1/0/2,345/2/1/0,348/2/1/2,456/2/2/0,467/2/3/1,478/2/3/2}{%
  \pgfmathsetmacro{\fiberlength}{5+\u-\w}%
  \pgfmathsetmacro{\dx}{0.34*\fiberlength}%
  \pgfmathsetmacro{\dy}{0.55*\fiberlength}%
  \coordinate (b\id) at (\u,\v,\w);%
  \coordinate (t\id) at ($(b\id)+(\dx cm,\dy cm)$);%
}

%================================================
% Fibers in the a-direction
%================================================
\foreach \id in {125,128,156,167,178,235,238,345,348,456,467,478}{%
  \draw[fiberthin] (b\id)--(t\id);%
}

% Highlight shortest and longest fibers
\draw[important] (b128)--(t128)
  node[pos=.50,left=4pt,font=\scriptsize,fill=white,inner sep=1pt] {$3$};

\draw[important] (b345)--(t345)
  node[pos=.55,right=4pt,font=\scriptsize,fill=white,inner sep=1pt] {$7$};

%================================================
% Graph of the bottom P_1
%================================================
\foreach \i/\j in {125/128,125/156,125/235,128/178,128/238,156/167,156/456,167/178,167/467,178/478,235/238,235/345,238/348,345/348,345/456,348/478,456/467,467/478}{%
  \draw[lowedge] (b\i)--(b\j);%
}

%================================================
% Graph of the top layer
%================================================
\foreach \i/\j in {125/128,125/156,125/235,128/178,128/238,156/167,156/456,167/178,167/467,178/478,235/238,235/345,238/348,345/348,345/456,348/478,456/467,467/478}{%
  \draw[topedge] (t\i)--(t\j);%
}

%================================================
% Vertex markers
%================================================
\foreach \id in {125,128,156,167,178,235,238,345,348,456,467,478}{%
  \fill[black] (b\id) circle (1.1pt);%
  \fill[blue!70!black] (t\id) circle (1.0pt);%
}

%================================================
% Bottom-layer vertex labels (all 12)
%================================================
\node[vlabel,below left=2pt]  at (b125) {$125$};
\node[vlabel,left=3pt]        at (b128) {$128$};
\node[vlabel,left=3pt]        at (b156) {$156$};
\node[vlabel,above left=1pt]  at ($(b167)+(-0.10,0.00)$) {$167$};
\node[vlabel,above left=2pt]  at (b178) {$178$};

\node[vlabel,below=3pt]       at (b235) {$235$};
\node[vlabel,above=4pt]       at ($(b238)+(0.02,0.00)$) {$238$};

\node[vlabel,below right=2pt] at (b345) {$345$};
\node[vlabel,above right=2pt] at ($(b348)+(0.08,0.00)$) {$348$};

\node[vlabel,below right=2pt] at (b456) {$456$};
\node[vlabel,right=3pt]       at ($(b467)+(0.06,0.00)$) {$467$};
\node[vlabel,above right=2pt] at (b478) {$478$};

%================================================
% Face labels for the lower P_1
%================================================
\node[flabel] at (-0.26,1.43,1.00) {$1$};
\node[flabel] at (0.40,-0.15,1.00) {$2$};
\node[flabel] at (1.45,0.34,1.03) {$3$};
\node[flabel] at (2.18,1.90,1.00) {$4$};
\node[flabel] at (1.00,1.02,-0.07) {$5$};
\node[flabel] at (1.08,2.48,0.30) {$6$};
\node[flabel] at (1.03,2.95,1.24) {$7$};
\node[flabel] at (1.02,1.55,2.10) {$8$};

%================================================
% Labels for the two layers
%================================================
\node[font=\small,text=black!75,anchor=east,fill=white,inner sep=1pt]
  at ($(b125)+(-0.60cm,-0.52cm)$) {$a=0$};

\node[font=\small,text=blue!65!black,anchor=west,fill=white,inner sep=1pt]
  at ($(t478)+(0.35cm,0.35cm)$) {$a=5+u-w$};


\node[font=\small,text=orange!80!black,anchor=west,fill=white,inner sep=1pt]
  at ($(t345)+(0.28cm,-0.35cm)$) {fiber length $=5+u-w$};

\end{tikzpicture}

\caption{
An affine projection of the four-dimensional lifted brick polytope $\mathcal P_2$.
The green part is the polytope $\mathcal P_1$. The red indices are facets of $P_1$. The blue part is a copy of $\mathcal P_1$. The orange segments represent the $a$-direction fibers.
}
\label{fig:P2-variable-fibers-refined}
\end{figure}


\subsection{Example 3}
We take $W=\widetilde A_2$,
$Q=s_1s_0s_1s_0s_2s_1s_2s_0s_1s_0s_1$ and
$\pi=\Dem(Q)=s_1s_0s_1s_2s_1s_0s_1$.
The length of $\pi$ is $7$ and $d=11-7=4$. Take
$I^+=\{1,2,5,8\}$. We have
\[
R(I^+)=(\alpha_1,\alpha_0,\alpha_1+\delta,\alpha_0+\delta),
\text{ where }\delta=\alpha_0+\alpha_1+\alpha_2.
\]
These roots are pairwise distinct up to sign, but
\[
\alpha_1-\alpha_0-(\alpha_1+\delta)+(\alpha_0+\delta)=0.
\]
This example is double root free but not root independent. 
Set
\[
M_Q=\Z A\oplus\Z B\oplus\Z C\oplus\Z D
\text{ and }
(A,B,C,D)=(e_1^{I^+},e_2^{I^+},e_5^{I^+},e_8^{I^+}).
\]
The projection from $M_Q$ to
$\Lambda_Q=\Z\alpha_0\oplus\Z\alpha_1\oplus\Z\alpha_2$ is
\begin{equation}\label{eq:int-P3-p}
p(a,b,c,d)=-a\alpha_1-b\alpha_0-c(\alpha_1+\delta)-d(\alpha_0+\delta).
\end{equation}
Its kernel is
\[
\ker p=\Z\kappa,\qquad \kappa=A-B-C+D.
\]

All $29$ facets and the coordinates of the lifted brick vectors after normalization are listed in Table 2.
\begin{table}[H]
\centering
\caption{$\SC_{\widetilde A_2}(Q,\pi)$: facets and lifted brick vectors.}
\label{tab:int-P3-vertices}
\medskip
\small
\begin{tabular}{@{}cc@{}}
\toprule
$I$ & $\wt B(I)=(a,b,c,d)$\\
\midrule
$\{1,2,5,8\}$ & $(0,0,0,0)$\\
$\{1,2,5,11\}$ & $(0,0,0,2)$\\
$\{1,2,8,9\}$ & $(0,0,2,0)$\\
$\{1,2,9,10\}$ & $(0,0,3,1)$\\
$\{1,2,10,11\}$ & $(0,0,3,2)$\\
$\{1,5,6,8\}$ & $(0,2,0,0)$\\
$\{1,5,6,11\}$ & $(0,2,0,2)$\\
$\{1,6,7,8\}$ & $(0,3,1,0)$\\
$\{1,6,7,11\}$ & $(0,3,1,2)$\\
$\{1,7,8,9\}$ & $(0,3,2,0)$\\
$\{1,7,9,10\}$ & $(0,3,3,1)$\\
$\{1,7,10,11\}$ & $(0,3,3,2)$\\
$\{2,3,5,8\}$ & $(1,0,0,0)$\\
$\{2,3,5,11\}$ & $(1,0,0,2)$\\
$\{2,3,8,9\}$ & $(1,0,2,0)$\\
$\{2,3,9,10\}$ & $(1,0,3,1)$\\
$\{2,3,10,11\}$ & $(1,0,3,2)$\\
$\{3,4,5,8\}$ & $(2,1,0,0)$\\
$\{3,4,5,11\}$ & $(2,1,0,2)$\\
$\{3,4,8,9\}$ & $(2,1,2,0)$\\
$\{3,4,9,10\}$ & $(2,1,3,1)$\\
$\{3,4,10,11\}$ & $(2,1,3,2)$\\
$\{4,5,6,8\}$ & $(2,2,0,0)$\\
$\{4,5,6,11\}$ & $(2,2,0,2)$\\
$\{4,6,7,8\}$ & $(2,3,1,0)$\\
$\{4,6,7,11\}$ & $(2,3,1,2)$\\
$\{4,7,8,9\}$ & $(2,3,2,0)$\\
$\{4,7,9,10\}$ & $(2,3,3,1)$\\
$\{4,7,10,11\}$ & $(2,3,3,2)$\\
\bottomrule
\end{tabular}
\end{table}

For the facet $I^+$, the brick vector is
\[
B(I^+)=4\omega_0+5\omega_1+2\omega_2-5\alpha_0-7\alpha_1-\alpha_2.
\]
The lifted brick polytope is
\begin{equation}\label{eq:int-P3}
\mathcal P_3=\{(a,b,c,d)\mid 0\le a\le2,\ 0\le b\le3,\ 0\le c\le3,\ 0\le d\le2,\
a-b\le1,\ b-c\le2,\ c-d\le2\}.
\end{equation}
It is a four-dimensional simple polytope with $29$ vertices and $11$ facets. Notice that the two distinct facets $I=\{1,7,8,9\}$ and $J=\{3,4,5,11\}$ satisfy $\wt B(J)-\wt B(I)=2\kappa \text{ and } B(J)=B(I).$ Thus the ordinary brick vectors do not distinguish these two facets.


